% filas del capítulo 10: símbolo & significado & unidad \\ (4.ª columna, macro, opcional; toda fila termina en \\)
$F$ & fracción de la operación designada como instrumento de cobertura & adimensional, $(0,1]$ \\
$h$ & razón de cobertura (unidades de instrumento por unidad de partida) & adimensional \\
$t_d$, $t_k$ & fecha de designación; fecha del cierre $k$ (con $t_0:=t_d$) & fecha \\
$\PV^{\mathrm{ins}}_t$, $\PV^{\mathrm{hip}}_t$ & valor presente del instrumento (operación completa) y del derivado hipotético & \MXN \\
$C^{\mathrm{ins}}_k$, $C^{\mathrm{hip}}_k$ & cambio acumulado de valor desde la designación (el del instrumento por $F$) & \MXN \\
$L^{\mathrm{ins}}_k$, $L^{\mathrm{hip}}_k$ & importe liquidado en el periodo $k$ (el del instrumento por $F$) & \MXN \\
$G^{\mathrm{ins}}_k$, $G^{\mathrm{hip}}_k$ & cambio acumulado más lo liquidado acumulado desde la designación & \MXN \\
$E_k$ & porción efectiva acumulada (regla del menor sobre $G$) & \MXN \\
$I_k$ & inefectividad del periodo, a resultados & \MXN \\
$q_k$, $Q_k$ & parte liquidada que coincide en el periodo y su acumulado & \MXN \\
$R_k$ & reciclaje del periodo, de la reserva a resultados & \MXN \\
$M_k$ & saldo de la reserva de cobertura (ORI) al cierre $k$ & \MXN \\
$a_k$ & ajuste de la partida cubierta en valor razonable & \MXN \\
$\eta_s$ & razón de la prueba prospectiva en el escenario $s$ & adimensional \\
$\kappa(x)$ & importe $x$ en centavos enteros (mitad alejándose de cero) & centavos \\
